% !TEX program = xelatex
\documentclass[11pt,a4paper]{article}

\usepackage[UTF8,fontset=none]{ctex}
\setCJKmainfont[BoldFont=Heiti SC,ItalicFont=Songti SC]{Songti SC}
\setCJKsansfont{Heiti SC}
\setCJKmonofont{Heiti SC}

\usepackage{amsmath,amssymb,amsfonts,mathtools,amsthm}
\usepackage{booktabs,array,longtable,tabularx}
\usepackage[margin=1.9cm]{geometry}
\usepackage[dvipsnames]{xcolor}
\usepackage[most]{tcolorbox}
\usepackage[hidelinks]{hyperref}
\usepackage{enumitem}
\usepackage{fancyhdr}

\setlength{\parindent}{0em}
\setlength{\parskip}{0.45em}
\setlist[itemize]{leftmargin=*,topsep=2pt,itemsep=1.5pt}
\setlist[enumerate]{leftmargin=*,topsep=2pt,itemsep=1.5pt}
\allowdisplaybreaks

\pagestyle{fancy}
\fancyhf{}
\lhead{\small Statistical Modelling and Inference}
\rhead{\small 2026 期末复习}
\cfoot{\thepage}

\hypersetup{
  pdftitle={Statistical Modelling and Inference 2026 Final Review},
  pdfauthor={Compiled from course notes, slides, and sample examinations}
}

\newcommand{\E}{\operatorname{E}}
\newcommand{\Var}{\operatorname{Var}}
\newcommand{\Cov}{\operatorname{Cov}}
\newcommand{\Bias}{\operatorname{Bias}}
\newcommand{\MSE}{\operatorname{MSE}}
\newcommand{\SE}{\operatorname{SE}}
\newcommand{\iid}{\overset{\mathrm{iid}}{\sim}}
\newcommand{\ind}{\perp\!\!\!\perp}
\newcommand{\1}{\mathbf{1}}
\newcommand{\bx}{\mathbf{x}}
\newcommand{\by}{\mathbf{y}}
\newcommand{\bX}{\mathbf{X}}
\newcommand{\bH}{\mathbf{H}}
\newcommand{\bI}{\mathbf{I}}
\newcommand{\be}{\mathbf{e}}
\newcommand{\bbeta}{\boldsymbol{\beta}}
\newcommand{\bvarepsilon}{\boldsymbol{\varepsilon}}

\newtcolorbox{scopebox}[1][]{
  colback=gray!5!white,colframe=gray!65!black,
  fonttitle=\bfseries,title=#1,breakable
}
\newtcolorbox{memorybox}[1][]{
  colback=orange!7!white,colframe=orange!75!black,
  fonttitle=\bfseries,title=#1,breakable
}
\newtcolorbox{conceptbox}[1][]{
  colback=blue!4!white,colframe=blue!60!black,
  fonttitle=\bfseries,title=#1,breakable
}
\newtcolorbox{methodbox}[1][]{
  colback=green!4!white,colframe=green!55!black,
  fonttitle=\bfseries,title=#1,breakable
}
\newtcolorbox{warnbox}[1][]{
  colback=red!4!white,colframe=red!60!black,
  fonttitle=\bfseries,title=#1,breakable
}
\newtcolorbox{examplebox}[1][]{
  colback=cyan!3!white,colframe=cyan!55!black,
  fonttitle=\bfseries,title=#1,breakable
}
\newtcolorbox{proofbox}[1][]{
  colback=violet!3!white,colframe=violet!55!black,
  fonttitle=\bfseries,title=#1,breakable
}

\title{\textbf{Statistical Modelling and Inference}\\[0.3em]
2026 期末复习讲义}
\author{依据黑板考试范围、课程讲义、课件、样卷与教师录音确认整理}
\date{June 2026}

\begin{document}
\maketitle

\begin{scopebox}[考试范围说明]
\textbf{本讲义纳入：}
\begin{enumerate}
  \item 无偏估计（unbiased estimation）与均方误差（MSE）；
  \item 估计与假设检验：单总体均值、两总体均值、总体方差未知时的均值推断、总体比例；
  \item 回归建模：简单线性回归、多元线性回归、最小二乘估计、残差及其性质；
  \item MLE 的推断：正态、二项、泊松、指数分布，包括 Fisher 信息与大样本近似推断。
\end{enumerate}

\textbf{明确排除：}ANOVA、R 编程、P-value、Power、Polynomial regression、nested-model comparison。

\textbf{教师录音确认：}期末共 5 道大题，每题可拆成若干小问；没有固定“几道证明、几道计算”的比例，但整体更偏 derivation/proof。Simple linear regression 仍需准备手算与基础推导；multiple linear regression 重点是理论推导与证明，不考给出大数据集后的繁重手算；polynomial regression 不纳入期末。

\textbf{第三张黑板图说明：}图中在 ``Unknown Variance'' 附近补写了两个统计量，但板书细节尚未由教师确认。本讲义当前按照课程材料中明确出现的内容，将 ``Unknown Variance'' 解释为“总体方差未知时，使用 $t$ 分布进行单总体或两总体均值推断”。不把尚未确认的板书解释当作考试结论。
\end{scopebox}

\tableofcontents
\newpage

% ============================================================
\section{全课程主线与做题决策}
% ============================================================

\subsection{从样本到结论}

统计推断的主线是
\[
  \text{总体参数未知}
  \quad \xleftarrow{\ \text{利用抽样分布}\ }\quad
  \text{样本统计量已观测}.
\]

\begin{center}
\begin{tabular}{lll}
\toprule
研究对象 & 总体参数（固定但未知） & 样本统计量（随样本变化）\\
\midrule
均值 & $\mu$ & $\bar X$\\
方差 & $\sigma^2$ & $S^2$\\
总体比例 & $p$ & $\hat p=X/n$\\
回归系数 & $\beta_j$ & $\hat\beta_j$\\
一般模型参数 & $\theta$ & $\hat\theta_{\mathrm{MLE}}$\\
\bottomrule
\end{tabular}
\end{center}

\begin{conceptbox}[Estimator 与 Estimate]
\textbf{估计量（estimator）}是样本的函数，如
\[
  \hat\theta=g(X_1,\ldots,X_n),
\]
因此在抽样之前是随机变量。\textbf{估计值（estimate）}是观察数据后代入得到的具体数值。证明无偏性、计算方差和 MSE 时讨论的是估计量；报告数据结果时给出的是估计值。
\end{conceptbox}

\subsection{推断方法选择总表}

\begin{center}
\renewcommand{\arraystretch}{1.25}
\begin{tabularx}{\textwidth}{p{2.4cm}p{3.5cm}p{3.6cm}X}
\toprule
目标参数 & 情形 & 核心标准误 & 参考分布\\
\midrule
$\mu$ & $\sigma$ 已知，或明确要求大样本正态近似
& $\sigma/\sqrt n$（大样本可按题意用 $s/\sqrt n$）
& 标准正态 $Z$\\
$\mu$ & $\sigma$ 未知且总体正态
& $s/\sqrt n$
& $t_{n-1}$\\
$\mu_1-\mu_2$ & 独立、大样本
& $\sqrt{s_1^2/n_1+s_2^2/n_2}$
& 标准正态 $Z$\\
$\mu_1-\mu_2$ & 独立、小样本、正态、等方差
& $s_p\sqrt{1/n_1+1/n_2}$
& $t_{n_1+n_2-2}$\\
$\mu_1-\mu_2$ & 独立、小样本、正态、不假设等方差
& $\sqrt{s_1^2/n_1+s_2^2/n_2}$
& Welch $t_\nu$\\
$\mu_d$ & 配对样本
& $s_d/\sqrt n$
& $t_{n-1}$\\
$p$ & 单总体比例、大样本
& CI 用 $\sqrt{\hat p(1-\hat p)/n}$；检验用 $\sqrt{p_0(1-p_0)/n}$
& 标准正态 $Z$\\
$p_1-p_2$ & 两独立比例、大样本
& CI 不合并；检验 $H_0:p_1=p_2$ 时合并
& 标准正态 $Z$\\
\bottomrule
\end{tabularx}
\end{center}

\begin{methodbox}[任何 CI / 检验题的第一分钟]
\begin{enumerate}
  \item 写清楚\textbf{参数}：究竟是 $\mu$、$\mu_1-\mu_2$、$\mu_d$、$p$ 还是 $p_1-p_2$。
  \item 判断\textbf{设计}：单样本、两个独立样本，还是配对样本。
  \item 判断\textbf{标准误与参考分布}：$Z$ 还是 $t$；方差是否未知；两样本是否等方差。
  \item 检查\textbf{假设条件}：随机性、独立性、正态性或大样本近似。
  \item 再代公式，最后用题目语境解释。
\end{enumerate}
\end{methodbox}

% ============================================================
\section{无偏估计、偏差与均方误差}
% ============================================================

\subsection{定义与基本性质}

\begin{memorybox}[必须记忆：Bias 与 MSE]
设 $\hat\theta$ 估计参数 $\theta$：
\[
  \boxed{\Bias(\hat\theta)=\E(\hat\theta)-\theta.}
\]
若 $\E(\hat\theta)=\theta$，则称 $\hat\theta$ 是 $\theta$ 的\textbf{无偏估计量}。
\[
  \boxed{\MSE(\hat\theta)=\E\bigl[(\hat\theta-\theta)^2\bigr].}
\]
\end{memorybox}

常用运算规则：
\begin{align*}
  \E(aX+b)&=a\E(X)+b,\\
  \Var(aX+b)&=a^2\Var(X),\\
  \Var(aX+bY)&=a^2\Var(X)+b^2\Var(Y)+2ab\Cov(X,Y).
\end{align*}
独立蕴含协方差为零，但协方差为零一般不蕴含独立。

\subsection{偏差--方差分解}

\begin{memorybox}[必须记忆：MSE 分解]
\[
  \boxed{
  \MSE(\hat\theta)
  =
  \Var(\hat\theta)+\bigl[\Bias(\hat\theta)\bigr]^2.
  }
\]
若 $\hat\theta$ 无偏，则 $\MSE(\hat\theta)=\Var(\hat\theta)$。
\end{memorybox}

\begin{proofbox}[完整证明]
把估计误差拆成随机波动与系统偏差：
\[
  \hat\theta-\theta
  =
  \bigl[\hat\theta-\E(\hat\theta)\bigr]
  +
  \bigl[\E(\hat\theta)-\theta\bigr].
\]
令 $B=\Bias(\hat\theta)=\E(\hat\theta)-\theta$，平方并取期望：
\begin{align*}
  \E[(\hat\theta-\theta)^2]
  &=
  \E\left[
  \bigl(\hat\theta-\E\hat\theta\bigr)^2
  +2B\bigl(\hat\theta-\E\hat\theta\bigr)
  +B^2
  \right]\\
  &=
  \Var(\hat\theta)
  +2B\,\underbrace{\E[\hat\theta-\E(\hat\theta)]}_{0}
  +B^2\\
  &=\Var(\hat\theta)+B^2.
\end{align*}
\end{proofbox}

\begin{conceptbox}[为什么 MSE 比“是否无偏”更全面]
无偏只说明长期平均位置正确，不保证每次估计都靠近真值。MSE 同时惩罚方差和偏差，因此一个轻微有偏但方差显著更小的估计量，可能比无偏估计量具有更小的 MSE。
\end{conceptbox}

\subsection{构造无偏估计量}

若
\[
  \E(\hat\theta)=a\theta+b,\qquad a\ne 0,
\]
则
\[
  \boxed{\hat\theta^*=\frac{\hat\theta-b}{a}}
\]
是无偏估计量，因为
\[
  \E(\hat\theta^*)=\frac{a\theta+b-b}{a}=\theta.
\]

\begin{examplebox}[样卷型：缩放后的样本均值]
设 $X_1,\ldots,X_n$ 的均值为 $\mu$、方差为 $\sigma^2$，并定义
\[
  T=\frac{n-1}{n}\bar X.
\]
则
\[
  \E(T)=\frac{n-1}{n}\mu,
  \qquad
  \Bias(T)=-\frac{\mu}{n},
\]
\[
  \Var(T)
  =
  \left(\frac{n-1}{n}\right)^2\frac{\sigma^2}{n}
  =
  \frac{(n-1)^2\sigma^2}{n^3},
\]
\[
  \MSE(T)
  =
  \frac{(n-1)^2\sigma^2}{n^3}
  +
  \frac{\mu^2}{n^2}.
\]
若令 $T^*=aT$ 并要求无偏，则
\[
  a\frac{n-1}{n}\mu=\mu
  \quad\Longrightarrow\quad
  a=\frac{n}{n-1},
\]
所以 $T^*=\bar X$。
\end{examplebox}

\subsection{合并两个无偏估计量}

设 $\hat\theta_1,\hat\theta_2$ 都无偏，定义
\[
  \hat\theta_a=a\hat\theta_1+(1-a)\hat\theta_2.
\]
由于权重和为 1，
\[
  \E(\hat\theta_a)=a\theta+(1-a)\theta=\theta.
\]
所以任意固定 $a$ 都保持无偏。

若
\[
  \Var(\hat\theta_1)=v_1,\quad
  \Var(\hat\theta_2)=v_2,\quad
  \Cov(\hat\theta_1,\hat\theta_2)=c,
\]
则
\[
  \Var(\hat\theta_a)
  =
  a^2v_1+(1-a)^2v_2+2a(1-a)c.
\]
因为无偏，最小化 MSE 等价于最小化方差。求导可得
\[
  \boxed{
  a^*=\frac{v_2-c}{v_1+v_2-2c}.
  }
\]
独立时 $c=0$：
\[
  \boxed{
  a^*=\frac{v_2}{v_1+v_2},
  \qquad
  1-a^*=\frac{v_1}{v_1+v_2}.
  }
\]
方差较小的估计量得到更大的权重。

\subsection{样本均值与样本方差}

若 $X_1,\ldots,X_n$ 独立同分布，$\E(X_i)=\mu$，$\Var(X_i)=\sigma^2$，则
\[
  \boxed{\E(\bar X)=\mu,\qquad \Var(\bar X)=\frac{\sigma^2}{n}.}
\]

定义
\[
  S_n^2=\frac{1}{n}\sum_{i=1}^n(X_i-\bar X)^2,
  \qquad
  S^2=\frac{1}{n-1}\sum_{i=1}^n(X_i-\bar X)^2.
\]
利用恒等式
\[
  \sum_{i=1}^n(X_i-\bar X)^2
  =
  \sum_{i=1}^n(X_i-\mu)^2-n(\bar X-\mu)^2,
\]
取期望：
\[
  \E\left[\sum_{i=1}^n(X_i-\bar X)^2\right]
  =
  n\sigma^2-n\frac{\sigma^2}{n}
  =
  (n-1)\sigma^2.
\]
因此
\[
  \boxed{\E(S_n^2)=\frac{n-1}{n}\sigma^2,\qquad \E(S^2)=\sigma^2.}
\]

\begin{warnbox}[常见陷阱]
\begin{itemize}
  \item 无偏样本方差用分母 $n-1$；正态方差的 MLE 通常用分母 $n$。
  \item $S^2$ 对 $\sigma^2$ 无偏，并不意味着 $S$ 对 $\sigma$ 无偏。
  \item 比较估计量不能只比较偏差；题目问“更好”时优先比较 MSE。
\end{itemize}
\end{warnbox}

% ============================================================
\section{Estimation 与 Testing 的统一框架}
% ============================================================

\subsection{点估计、区间估计与假设检验}

\textbf{点估计}给一个最合理的数值；\textbf{置信区间（CI）}给出与数据相容的一段参数范围；\textbf{假设检验}考察某个指定参数值是否与数据相容。

大多数双侧 CI 都具有结构
\[
  \boxed{
  \text{point estimate}
  \ \pm\
  \text{critical value}\times \text{standard error}.
  }
\]

\begin{conceptbox}[置信水平的正确解释]
一个具体算出的 95\% CI 要么覆盖真参数，要么不覆盖；在频率学派中参数不是随机的。95\% 的含义是：若无限次重复抽样并用同一方法构造区间，约 95\% 的区间会覆盖真参数。
\end{conceptbox}

\subsection{假设检验的结构}

原假设必须包含参数空间中的\textbf{边界值}，可以是简单原假设
\[
  H_0:\theta=\theta_0.
\]
也可以是单侧复合原假设，例如
\[
  H_0:\theta\leq\theta_0
  \quad\text{vs.}\quad
  H_a:\theta>\theta_0.
\]
在后一种情形中，临界值法通常在原假设边界 $\theta=\theta_0$ 处计算检验统计量的参考分布。不能只因为观测估计值落在备择方向，就直接拒绝原假设。

备择假设决定尾部方向：
\[
  H_a:\theta>\theta_0,\qquad
  H_a:\theta<\theta_0,\qquad
  H_a:\theta\ne\theta_0.
\]

标准化统计量的基本结构是
\[
  \boxed{
  \text{test statistic}
  =
  \frac{\text{estimate}-\text{null value}}{\text{standard error under }H_0}.
  }
\]
它衡量观测估计值距离原假设值有多少个标准误。

\begin{conceptbox}[第一类错误与第二类错误]
\begin{itemize}
  \item \textbf{第一类错误（Type I error）}：$H_0$ 真实却拒绝 $H_0$。检验设计使其概率不超过显著性水平 $\alpha$；在边界处通常等于 $\alpha$。
  \item \textbf{第二类错误（Type II error）}：$H_0$ 错误却不能拒绝 $H_0$。其概率通常记为 $\beta$，会随真实参数值、样本量和拒绝域变化。
\end{itemize}
“不能拒绝 $H_0$”不等于证明 $H_0$ 为真；它只表示现有样本没有提供足够强的反对证据。
\end{conceptbox}

\begin{memorybox}[临界值法：本讲义唯一采用的检验决策方法]
若检验统计量在 $H_0$ 下服从标准正态：
\[
\begin{array}{lll}
H_a:\theta>\theta_0 &\Rightarrow& Z_{\mathrm{obs}}>z_\alpha\ \text{时拒绝 }H_0,\\
H_a:\theta<\theta_0 &\Rightarrow& Z_{\mathrm{obs}}<-z_\alpha\ \text{时拒绝 }H_0,\\
H_a:\theta\ne\theta_0 &\Rightarrow& |Z_{\mathrm{obs}}|>z_{\alpha/2}\ \text{时拒绝 }H_0.
\end{array}
\]
$t$ 检验将 $z$ 临界值替换为相应自由度的 $t$ 临界值。
\end{memorybox}

\begin{center}
\begin{tabular}{cccc}
\toprule
置信水平 & $\alpha$ & 双侧临界值 $z_{\alpha/2}$ & 单侧临界值 $z_\alpha$\\
\midrule
90\% & 0.10 & 1.645 & 1.282\\
95\% & 0.05 & 1.960 & 1.645\\
99\% & 0.01 & 2.576 & 2.326\\
\bottomrule
\end{tabular}
\end{center}

\begin{methodbox}[假设检验标准答题模板]
\begin{enumerate}
  \item 定义参数，并写出 $H_0$ 与 $H_a$；
  \item 写出显著性水平 $\alpha$；
  \item 说明适用条件与所用参考分布；
  \item 写公式并计算观测检验统计量；
  \item 写出临界值和拒绝域；
  \item 比较并作出“拒绝 $H_0$”或“不能拒绝 $H_0$”的决定；
  \item 用题目语境写结论，不写“接受 $H_0$”。
\end{enumerate}
\end{methodbox}

\subsection{CI 与双侧检验的关系}

在使用同一标准误与同一参考分布时，显著性水平为 $\alpha$ 的双侧检验
\[
  H_0:\theta=\theta_0
  \quad\text{vs.}\quad
  H_a:\theta\ne\theta_0
\]
与 $(1-\alpha)100\%$ CI 具有如下对应关系：
\[
  \boxed{
  \theta_0\notin \text{CI}
  \Longleftrightarrow
  \text{拒绝 }H_0.
  }
\]

% ============================================================
\section{单总体均值：已知与未知方差}
% ============================================================

\subsection{抽样分布}

若 $X_1,\ldots,X_n$ 独立同分布，均值为 $\mu$、方差为 $\sigma^2$：
\[
  \E(\bar X)=\mu,\qquad
  \SE(\bar X)=\frac{\sigma}{\sqrt n}.
\]
若总体正态，则 $\bar X$ 精确正态；若总体不正态但样本量足够大，则由中心极限定理近似正态。

\subsection{\texorpdfstring{总体标准差已知：$Z$ 推断}{总体标准差已知：Z 推断}}

\begin{memorybox}[单总体均值：$\sigma$ 已知]
\[
  Z=\frac{\bar X-\mu}{\sigma/\sqrt n}\sim N(0,1).
\]
$(1-\alpha)100\%$ CI：
\[
  \boxed{\bar x\pm z_{\alpha/2}\frac{\sigma}{\sqrt n}.}
\]
检验 $H_0:\mu=\mu_0$：
\[
  \boxed{Z_{\mathrm{obs}}=\frac{\bar x-\mu_0}{\sigma/\sqrt n}.}
\]
\end{memorybox}

\begin{examplebox}[单均值 $Z$ 区间与检验]
已知 $n=64,\bar x=33,\sigma=16$。95\% CI 为
\[
  33\pm1.96\frac{16}{8}
  =
  33\pm3.92
  =
  (29.08,36.92).
\]
检验 $H_0:\mu=30$ 对 $H_a:\mu>30$，$\alpha=0.05$：
\[
  Z_{\mathrm{obs}}=\frac{33-30}{16/\sqrt{64}}=1.50.
\]
右尾临界值为 $z_{0.05}=1.645$。因为 $1.50<1.645$，不能拒绝 $H_0$；数据不足以支持总体平均购物时间超过 30 分钟。
\end{examplebox}

\subsection{\texorpdfstring{总体方差未知：$t$ 推断}{总体方差未知：t 推断}}

当 $\sigma$ 未知，用样本标准差 $S$ 替代。若总体正态，则
\[
  \boxed{
  T=\frac{\bar X-\mu}{S/\sqrt n}\sim t_{n-1}.
  }
\]
$t$ 分布比标准正态尾部更厚，反映了用 $S$ 估计 $\sigma$ 带来的额外不确定性；自由度增加时，$t$ 分布趋近标准正态。

\begin{memorybox}[单总体均值：$\sigma$ 未知]
$(1-\alpha)100\%$ CI：
\[
  \boxed{\bar x\pm t_{\alpha/2,n-1}\frac{s}{\sqrt n}.}
\]
检验 $H_0:\mu=\mu_0$：
\[
  \boxed{
  t_{\mathrm{obs}}
  =
  \frac{\bar x-\mu_0}{s/\sqrt n},
  \qquad \nu=n-1.
  }
\]
\end{memorybox}

\begin{examplebox}[单均值 $t$ 区间]
某总体近似正态，$n=24,\bar x=98.25,s=0.73$。若
$t_{0.05,23}=1.714$，则 90\% CI 为
\[
  98.25\pm1.714\frac{0.73}{\sqrt{24}}
  =
  98.25\pm0.255
  =
  (97.995,98.505).
\]
\end{examplebox}

\begin{warnbox}[Unknown Variance 的关键含义]
\begin{itemize}
  \item “方差未知”不意味着不能推断；它意味着用 $S$ 替代 $\sigma$，并使用 $t$ 参考分布。
  \item 小样本 $t$ 推断要求总体近似正态；样本量很大时方法对轻度非正态更稳健。
  \item 自由度是 $n-1$，因为样本方差计算中先估计了一个均值。
\end{itemize}
\end{warnbox}

% ============================================================
\section{两总体均值的估计与检验}
% ============================================================

\subsection{先判断：独立还是配对}

\begin{conceptbox}[独立样本 vs 配对样本]
\textbf{独立样本：}两组由不同个体组成，例如治疗组与对照组。\\
\textbf{配对样本：}同一个体前后测量，或按特征一一匹配。配对设计的分析单位是差值
\[
  D_i=X_{1i}-X_{2i},
\]
因此本质上是对差值做单样本推断。
\end{conceptbox}

\subsection{两个独立大样本}

若两组独立且样本量足够大，
\[
  \SE(\bar X_1-\bar X_2)
  =
  \sqrt{\frac{s_1^2}{n_1}+\frac{s_2^2}{n_2}}.
\]

\begin{memorybox}[两个独立大样本均值]
$(1-\alpha)100\%$ CI：
\[
  \boxed{
  (\bar x_1-\bar x_2)
  \pm
  z_{\alpha/2}
  \sqrt{\frac{s_1^2}{n_1}+\frac{s_2^2}{n_2}}.
  }
\]
检验 $H_0:\mu_1-\mu_2=\Delta_0$：
\[
  \boxed{
  Z_{\mathrm{obs}}
  =
  \frac{(\bar x_1-\bar x_2)-\Delta_0}
  {\sqrt{s_1^2/n_1+s_2^2/n_2}}.
  }
\]
\end{memorybox}

\begin{examplebox}[两个独立大样本]
组 A：$n_1=36,\bar x_1=18.4,s_1=4.8$；组 B：
$n_2=49,\bar x_2=15.9,s_2=5.6$。
\[
  \widehat{\mu_A-\mu_B}=2.5,
\]
\[
  \SE
  =
  \sqrt{\frac{4.8^2}{36}+\frac{5.6^2}{49}}
  =
  \sqrt{0.64+0.64}
  =
  1.131.
\]
95\% CI：
\[
  2.5\pm1.96(1.131)
  =
  (0.283,4.717).
\]
检验 $H_0:\mu_A-\mu_B=0$ 对 $H_a:\mu_A-\mu_B>0$：
\[
  Z_{\mathrm{obs}}=\frac{2.5}{1.131}=2.21>1.645.
\]
拒绝 $H_0$；数据支持组 A 的总体平均学习时间更长。
\end{examplebox}

\subsection{\texorpdfstring{小样本、正态、等方差：Pooled $t$}{小样本、正态、等方差：Pooled t}}

若两个总体正态、相互独立，并可假设
\[
  \sigma_1^2=\sigma_2^2=\sigma^2,
\]
则把两个样本方差按自由度加权合并：
\[
  \boxed{
  s_p^2
  =
  \frac{(n_1-1)s_1^2+(n_2-1)s_2^2}{n_1+n_2-2}.
  }
\]

\begin{memorybox}[两个独立小样本：等方差]
\[
  \boxed{
  t_{\mathrm{obs}}
  =
  \frac{(\bar x_1-\bar x_2)-\Delta_0}
  {s_p\sqrt{1/n_1+1/n_2}},
  \qquad
  \nu=n_1+n_2-2.
  }
\]
\[
  \boxed{
  (\bar x_1-\bar x_2)
  \pm
  t_{\alpha/2,\nu}
  s_p\sqrt{\frac1{n_1}+\frac1{n_2}}.
  }
\]
\end{memorybox}

\subsection{\texorpdfstring{扩展：小样本、不假设等方差的 Welch $t$}{扩展：小样本、不假设等方差的 Welch t}}

Welch 推断是处理两个独立小样本且不假设等方差的常用方法，但课程材料主要强调 pooled $t$。若试题未要求 Welch，先依据题目给出的假设与课堂约定选法。Welch 方法为：
\[
  \SE
  =
  \sqrt{\frac{s_1^2}{n_1}+\frac{s_2^2}{n_2}},
\qquad
  t_{\mathrm{obs}}
  =
  \frac{(\bar x_1-\bar x_2)-\Delta_0}{\SE}.
\]
近似自由度为
\[
  \boxed{
  \nu
  \approx
  \frac{\left(s_1^2/n_1+s_2^2/n_2\right)^2}
  {\dfrac{(s_1^2/n_1)^2}{n_1-1}
  +
  \dfrac{(s_2^2/n_2)^2}{n_2-1}}.
  }
\]
若试题明确要求 pooled $t$ 或明确给出等方差，则使用上一节方法。

\subsection{配对样本}

定义差值
\[
  D_i=X_{1i}-X_{2i},
\qquad i=1,\ldots,n.
\]
计算差值样本均值 $\bar d$ 与标准差 $s_d$，再做单总体 $t$ 推断。

\begin{memorybox}[配对 $t$ 推断]
\[
  \boxed{
  t_{\mathrm{obs}}
  =
  \frac{\bar d-\Delta_0}{s_d/\sqrt n},
  \qquad \nu=n-1.
  }
\]
\[
  \boxed{
  \bar d\pm t_{\alpha/2,n-1}\frac{s_d}{\sqrt n}.
  }
\]
\end{memorybox}

\begin{warnbox}[两总体均值高频错误]
\begin{itemize}
  \item 配对数据不能当成两个独立样本；必须先求每一对的差值。
  \item 独立大样本标准误不能写成 $s_1/\sqrt{n_1}+s_2/\sqrt{n_2}$；方差先相加，再开方。
  \item Pooled $t$ 需要等方差假设；Welch $t$ 不合并方差。
  \item 参数顺序必须与样本差的顺序一致。例如估计 $\mu_A-\mu_B$ 就应计算 $\bar x_A-\bar x_B$。
\end{itemize}
\end{warnbox}

% ============================================================
\section{总体比例的估计与检验}
% ============================================================

\subsection{单总体比例}

若 $X$ 是 $n$ 次独立 Bernoulli 试验中的成功次数，则
\[
  X\sim\operatorname{Binomial}(n,p),
  \qquad
  \hat p=\frac{X}{n}.
\]
\[
  \E(\hat p)=p,\qquad
  \Var(\hat p)=\frac{p(1-p)}{n}.
\]

\begin{memorybox}[单总体比例：CI 与检验分母不同]
CI 使用样本比例估计标准误：
\[
  \boxed{
  \hat p
  \pm
  z_{\alpha/2}
  \sqrt{\frac{\hat p(1-\hat p)}{n}}.
  }
\]
检验 $H_0:p=p_0$ 时，标准误在 $H_0$ 下计算：
\[
  \boxed{
  Z_{\mathrm{obs}}
  =
  \frac{\hat p-p_0}
  {\sqrt{p_0(1-p_0)/n}}.
  }
\]
\end{memorybox}

正态近似通常要求成功与失败的期望次数都足够大。常用检查为
\[
  n\hat p\ge 5,\quad n(1-\hat p)\ge 5
\]
用于 CI；检验时用 $np_0$ 与 $n(1-p_0)$ 检查。

\begin{examplebox}[单总体比例]
在 $n=400$ 人中有 $x=100$ 人具有某特征，因此 $\hat p=0.25$。
检验 $H_0:p=0.20$ 对 $H_a:p\ne0.20$，$\alpha=0.05$：
\[
  Z_{\mathrm{obs}}
  =
  \frac{0.25-0.20}{\sqrt{0.20(0.80)/400}}
  =
  \frac{0.05}{0.02}
  =
  2.50.
\]
因为 $|2.50|>1.96$，拒绝 $H_0$；数据支持总体比例不同于 20\%。
\end{examplebox}

\subsection{两个独立总体比例}

设
\[
  \hat p_1=\frac{x_1}{n_1},
  \qquad
  \hat p_2=\frac{x_2}{n_2}.
\]

\begin{memorybox}[两总体比例：CI 不合并]
\[
  \boxed{
  (\hat p_1-\hat p_2)
  \pm
  z_{\alpha/2}
  \sqrt{
  \frac{\hat p_1(1-\hat p_1)}{n_1}
  +
  \frac{\hat p_2(1-\hat p_2)}{n_2}
  }.
  }
\]
\end{memorybox}

检验
\[
  H_0:p_1=p_2
\]
时，原假设认为两组共享同一个比例，因此先计算合并比例
\[
  \hat p_{\mathrm{pool}}
  =
  \frac{x_1+x_2}{n_1+n_2},
\]
再使用
\[
  \boxed{
  Z_{\mathrm{obs}}
  =
  \frac{\hat p_1-\hat p_2}
  {\sqrt{
  \hat p_{\mathrm{pool}}(1-\hat p_{\mathrm{pool}})
  \left(1/n_1+1/n_2\right)
  }}.
  }
\]

\begin{warnbox}[比例推断最常见陷阱]
\begin{itemize}
  \item 单比例 CI 分母用 $\hat p$，单比例检验分母用 $p_0$。
  \item 两比例 CI 不合并；检验 $H_0:p_1=p_2$ 时才合并。
  \item 两比例 CI 的第二项必须是
  $\hat p_2(1-\hat p_2)/n_2$；原课程 PDF 的该处存在下标笔误。
  \item 区间端点若超出 $[0,1]$，说明简单正态近似可能不理想。
\end{itemize}
\end{warnbox}

% ============================================================
\section{回归建模总览}
% ============================================================

\subsection{回归的目标}

回归分析用于描述和预测响应变量 $Y$ 如何随预测变量变化。必须区分：
\[
  \text{条件均值 } \E(Y\mid X=x)
  \quad\text{与}\quad
  \text{单个观测 }Y.
\]
即使条件均值是一条确定的曲线，单个观测仍因误差项而随机波动。

\subsection{误差假设}

简单线性回归的总体模型为
\[
  Y_i=\beta_0+\beta_1x_i+\varepsilon_i.
\]
经典假设可写为：
\[
  \E(\varepsilon_i\mid x_i)=0,
  \qquad
  \Var(\varepsilon_i\mid x_i)=\sigma^2,
  \qquad
  \varepsilon_i\ind\varepsilon_j\ (i\ne j).
\]
若还假定
\[
  \varepsilon_i\sim N(0,\sigma^2),
\]
则可得到精确的小样本 $t$ 推断。正态性主要服务于推断，不是最小二乘系数存在的必要条件。

\begin{conceptbox}[模型假设与残差图]
\begin{itemize}
  \item 线性或函数形式合理：残差对拟合值图不应出现系统曲线。
  \item 同方差：残差散布宽度应大致恒定，不应呈漏斗形。
  \item 独立：依赖研究设计与时间顺序，不能只靠普通散点图保证。
  \item 正态性：误差应近似正态；残差只能作为误差的观测代理。
\end{itemize}
\end{conceptbox}

% ============================================================
\section{简单线性回归与最小二乘估计}
% ============================================================

\subsection{模型与系数解释}

\[
  Y_i=\beta_0+\beta_1x_i+\varepsilon_i,
  \qquad
  \E(Y_i\mid x_i)=\beta_0+\beta_1x_i.
\]
\begin{itemize}
  \item $\beta_0$：$x=0$ 时的总体平均响应；若 $0$ 不在数据合理范围内，它可能只具有数学意义。
  \item $\beta_1$：$x$ 每增加一个单位，条件均值平均变化 $\beta_1$ 个单位。
\end{itemize}

\subsection{最小二乘准则与正规方程}

拟合直线
\[
  \hat y_i=b_0+b_1x_i
\]
使残差平方和最小：
\[
  SSE(b_0,b_1)
  =
  \sum_{i=1}^n(y_i-b_0-b_1x_i)^2.
\]
对 $b_0,b_1$ 求偏导并令其为零：
\[
  \frac{\partial SSE}{\partial b_0}
  =
  -2\sum_{i=1}^n(y_i-b_0-b_1x_i)=0,
\]
\[
  \frac{\partial SSE}{\partial b_1}
  =
  -2\sum_{i=1}^nx_i(y_i-b_0-b_1x_i)=0.
\]
这给出正规方程，也立刻给出残差的两个正交性质：
\[
  \sum e_i=0,
  \qquad
  \sum x_ie_i=0.
\]

定义
\[
  S_{xx}=\sum_{i=1}^n(x_i-\bar x)^2,
\quad
  S_{xy}=\sum_{i=1}^n(x_i-\bar x)(y_i-\bar y),
\quad
  S_{yy}=\sum_{i=1}^n(y_i-\bar y)^2.
\]

\begin{memorybox}[简单线性回归 OLS]
\[
  \boxed{\hat\beta_1=b_1=\frac{S_{xy}}{S_{xx}}.}
\]
\[
  \boxed{\hat\beta_0=b_0=\bar y-b_1\bar x.}
\]
\[
  \boxed{\hat y=b_0+b_1x.}
\]
\end{memorybox}

\subsection{最小二乘直线的几何性质}

\begin{itemize}
  \item 拟合直线必经过 $(\bar x,\bar y)$：
  \[
    b_0+b_1\bar x=\bar y.
  \]
  \item 若模型含截距，残差均值为零：
  \[
    \bar e=0.
  \]
  \item 残差与预测变量正交：
  \[
    \sum x_ie_i=0.
  \]
  \item 残差与拟合值正交：
  \[
    \sum \hat y_ie_i=0.
  \]
\end{itemize}

\subsection{平方和与误差方差}

\[
  SST=S_{yy}=\sum(y_i-\bar y)^2,
\]
\[
  SSE=\sum(y_i-\hat y_i)^2
  =
  S_{yy}-\frac{S_{xy}^2}{S_{xx}}
  =
  S_{yy}-\hat\beta_1S_{xy}.
\]
简单线性回归估计了两个系数，故误差方差的无偏估计为
\[
  \boxed{
  s_e^2=\frac{SSE}{n-2},
  \qquad
  s_e=\sqrt{\frac{SSE}{n-2}}.
  }
\]

\subsection{平方和分解、相关系数与决定系数}

在含截距的简单线性回归中，总变异可分解为模型解释的变异与未解释的残差变异：
\[
  \boxed{
  SST=SSR+SSE,
  }
\]
其中
\[
  SST=\sum_{i=1}^n(y_i-\bar y)^2,
  \qquad
  SSR=\sum_{i=1}^n(\hat y_i-\bar y)^2,
  \qquad
  SSE=\sum_{i=1}^n(y_i-\hat y_i)^2.
\]
决定系数为
\[
  \boxed{
  R^2=\frac{SSR}{SST}=1-\frac{SSE}{SST}.
  }
\]
$R^2$ 表示样本中 $Y$ 的总变异有多大比例被拟合模型解释。它不证明因果关系，也不能单独保证模型形式正确。

样本相关系数为
\[
  \boxed{
  r
  =
  \frac{S_{xy}}{\sqrt{S_{xx}S_{yy}}},
  \qquad -1\le r\le1.
  }
\]
在\textbf{含截距的简单线性回归}中，
\[
  \boxed{R^2=r^2.}
\]
但在多元回归中，不再有某一个简单相关系数满足该等式。

\subsection{Week 9：OLS 估计量性质的标量证明}

令
\[
  a_i=\frac{x_i-\bar x}{S_{xx}},
\qquad
  \hat\beta_1=\sum_{i=1}^na_iY_i.
\]
又因为
\[
  \hat\beta_0=\bar Y-\bar x\hat\beta_1,
\]
可写成
\[
  \hat\beta_0=\sum_{i=1}^nb_iY_i,
  \qquad
  b_i=\frac1n-\bar x a_i.
\]
权重满足
\[
  \sum a_i=0,
  \qquad
  \sum a_ix_i=1,
  \qquad
  \sum a_i^2=\frac1{S_{xx}}.
\]
同时
\[
  \sum b_i=1,
  \qquad
  \sum b_ix_i=0,
\]
\[
  \sum b_i^2=\frac1n+\frac{\bar x^2}{S_{xx}},
  \qquad
  \sum a_ib_i=-\frac{\bar x}{S_{xx}}.
\]

\begin{conceptbox}[老师手写证明的入口]
手写备课笔记的核心技巧是把估计量写成 $Y_i$ 的线性组合：
\[
  \hat\beta_1=\sum a_iY_i,
  \qquad
  \hat\beta_0=\sum b_iY_i.
\]
之后所有期望、方差和协方差证明，都只是在使用
\[
  \E(Y_i)=\beta_0+\beta_1x_i,\qquad
  \Var(Y_i)=\sigma^2,\qquad
  Y_i\text{ 相互独立}.
\]
\end{conceptbox}

\begin{proofbox}[$\hat\beta_1$ 无偏性的完整证明]
\begin{align*}
  \E(\hat\beta_1)
  &=
  \sum_{i=1}^na_i\E(Y_i)\\
  &=
  \sum_{i=1}^na_i(\beta_0+\beta_1x_i)\\
  &=
  \beta_0\sum_{i=1}^na_i
  +
  \beta_1\sum_{i=1}^na_ix_i\\
  &=0+\beta_1=\beta_1.
\end{align*}
\end{proofbox}

\begin{proofbox}[$\hat\beta_0$ 无偏性的完整证明]
\begin{align*}
  \E(\hat\beta_0)
  &=
  \sum_{i=1}^nb_i\E(Y_i)\\
  &=
  \sum_{i=1}^nb_i(\beta_0+\beta_1x_i)\\
  &=
  \beta_0\sum_{i=1}^nb_i
  +
  \beta_1\sum_{i=1}^nb_ix_i\\
  &=\beta_0+0=\beta_0.
\end{align*}
\end{proofbox}

由独立与同方差：
\[
  \Var(\hat\beta_1)
  =
  \sum a_i^2\sigma^2
  =
  \boxed{\frac{\sigma^2}{S_{xx}}}.
\]
\[
  \Var(\hat\beta_0)
  =
  \sum b_i^2\sigma^2
  =
  \sigma^2\left(\frac1n+\frac{\bar x^2}{S_{xx}}\right).
\]
以及
\[
  \Cov(\hat\beta_0,\hat\beta_1)
  =
  \Cov\left(\sum b_iY_i,\sum a_iY_i\right)
  =
  \sigma^2\sum a_ib_i
  =
  -\frac{\sigma^2\bar x}{S_{xx}}.
\]
因此
\[
  \boxed{
  \Var(\hat\beta_0)
  =
  \sigma^2\left(\frac1n+\frac{\bar x^2}{S_{xx}}\right),
  }
\]
\[
  \boxed{
  \Cov(\hat\beta_0,\hat\beta_1)
  =
  -\frac{\sigma^2\bar x}{S_{xx}}.
  }
\]
若 $\bar x>0$，斜率与截距估计负相关；斜率更陡时，为保持直线经过 $(\bar x,\bar y)$，截距倾向降低。

\subsection{拟合均值与单个未来观测}

在 $x_0$ 处的拟合均值为
\[
  \hat Y_0=\hat\beta_0+\hat\beta_1x_0.
\]
它是 $\E(Y\mid X=x_0)=\beta_0+\beta_1x_0$ 的无偏估计：
\[
  \E(\hat Y_0)=\beta_0+\beta_1x_0.
\]
方差为
\[
  \boxed{
  \Var(\hat Y_0)
  =
  \sigma^2
  \left[
  \frac1n+\frac{(x_0-\bar x)^2}{S_{xx}}
  \right].
  }
\]
它在 $x_0=\bar x$ 处最小，离样本中心越远越大。

均值响应的 CI：
\[
  \boxed{
  \hat y_0
  \pm
  t_{\alpha/2,n-2}s_e
  \sqrt{\frac1n+\frac{(x_0-\bar x)^2}{S_{xx}}}.
  }
\]
单个未来观测还包含新误差项，因此预测区间更宽：
\[
  \boxed{
  \hat y_0
  \pm
  t_{\alpha/2,n-2}s_e
  \sqrt{1+\frac1n+\frac{(x_0-\bar x)^2}{S_{xx}}}.
  }
\]

\subsection{斜率的区间与检验}

用 $s_e^2$ 替代 $\sigma^2$：
\[
  \SE(\hat\beta_1)=\frac{s_e}{\sqrt{S_{xx}}}.
\]
斜率 CI：
\[
  \boxed{
  \hat\beta_1
  \pm
  t_{\alpha/2,n-2}\frac{s_e}{\sqrt{S_{xx}}}.
  }
\]
检验 $H_0:\beta_1=\beta_{1,0}$：
\[
  \boxed{
  t_{\mathrm{obs}}
  =
  \frac{\hat\beta_1-\beta_{1,0}}
  {s_e/\sqrt{S_{xx}}},
  \qquad \nu=n-2.
  }
\]

% ============================================================
\section{残差及其性质}
% ============================================================

\subsection{误差与残差的区别}

\[
  \varepsilon_i=Y_i-\E(Y_i\mid x_i)
\]
是不可直接观测的真实误差；而
\[
  e_i=y_i-\hat y_i
\]
是拟合模型后可计算的残差。残差用于诊断误差假设，但残差不是误差本身。

\subsection{简单线性回归残差性质}

\begin{memorybox}[必须掌握：含截距 OLS 残差]
\[
  \boxed{\sum_{i=1}^ne_i=0,}
\qquad
  \boxed{\sum_{i=1}^nx_ie_i=0,}
\qquad
  \boxed{\sum_{i=1}^n\hat y_ie_i=0.}
\]
在模型正确时：
\[
  \boxed{\E(e_i)=0.}
\]
\end{memorybox}

简单线性回归中，第 $i$ 个观测的杠杆值为
\[
  h_{ii}
  =
  \frac1n+\frac{(x_i-\bar x)^2}{S_{xx}}.
\]
残差方差为
\[
  \boxed{
  \Var(e_i)
  =
  \sigma^2(1-h_{ii})
  =
  \sigma^2
  \left[
  1-\frac1n-\frac{(x_i-\bar x)^2}{S_{xx}}
  \right].
  }
\]

\begin{proofbox}[Week 9：残差期望与方差的标量证明]
在简单线性回归中
\[
  e_i
  =
  Y_i-\hat\beta_0-\hat\beta_1x_i
  =
  Y_i-\bar Y-(x_i-\bar x)\hat\beta_1.
\]
因此
\[
\begin{aligned}
  \E(e_i)
  &=
  \E(Y_i)-\E(\bar Y)-(x_i-\bar x)\E(\hat\beta_1)\\
  &=
  (\beta_0+\beta_1x_i)-(\beta_0+\beta_1\bar x)
  -(x_i-\bar x)\beta_1
  =
  0.
\end{aligned}
\]
令 $d_i=x_i-\bar x$。由于
\[
  \Var(Y_i-\bar Y)=\sigma^2\left(1-\frac1n\right),
  \qquad
  \Cov(Y_i-\bar Y,\hat\beta_1)=\frac{d_i\sigma^2}{S_{xx}},
\]
并且
\[
  \Var(\hat\beta_1)=\frac{\sigma^2}{S_{xx}},
\]
所以
\[
\begin{aligned}
  \Var(e_i)
  &=
  \Var(Y_i-\bar Y-d_i\hat\beta_1)\\
  &=
  \sigma^2\left(1-\frac1n\right)
  +
  d_i^2\frac{\sigma^2}{S_{xx}}
  -
  2d_i\frac{d_i\sigma^2}{S_{xx}}\\
  &=
  \sigma^2\left[
    1-\frac1n-\frac{(x_i-\bar x)^2}{S_{xx}}
  \right].
\end{aligned}
\]
\end{proofbox}

\begin{conceptbox}[为什么残差方差小于误差方差]
拟合值使用了同一个观测的信息，因此残差不是未经处理的原始误差。高杠杆观测对自己的拟合值影响更大，故其原始残差方差反而更小。这也是诊断时常使用标准化或学生化残差的原因。
\end{conceptbox}

\subsection{残差诊断的模式识别}

\begin{center}
\begin{tabularx}{\textwidth}{p{3.7cm}X}
\toprule
残差图现象 & 可能问题\\
\midrule
围绕 0 随机散布、宽度大致恒定 & 线性与同方差假设较合理\\
明显弯曲或波浪形 & 函数形式错误，可能需要重新指定模型或作变量变换\\
漏斗形，散布随拟合值增大 & 异方差\\
孤立的大残差 & 可能的异常响应值\\
远离数据中心且杠杆高 & 可能的高杠杆点；不一定有大残差\\
时间顺序上连续同号或周期模式 & 独立性可能失效\\
\bottomrule
\end{tabularx}
\end{center}

% ============================================================
\section{多元线性回归与矩阵最小二乘}
% ============================================================

\subsection{模型与设计矩阵}

含 $k$ 个预测变量的模型为
\[
  Y_i=\beta_0+\beta_1x_{i1}+\cdots+\beta_kx_{ik}+\varepsilon_i.
\]
矩阵形式：
\[
  \boxed{\by=\bX\bbeta+\bvarepsilon.}
\]
其中 $\bX$ 的第一列通常为全 1 截距列，维数为
\[
  \bX:n\times p,
  \qquad p=k+1.
\]

\begin{conceptbox}[多元回归系数的解释]
$\beta_j$ 表示：\textbf{保持其他预测变量不变}时，$x_j$ 增加一个单位，条件平均响应变化 $\beta_j$ 个单位。它不是忽略其他变量后 $x_j$ 与 $Y$ 的简单关联。
\end{conceptbox}

\subsection{矩阵最小二乘推导}

最小化
\[
  SSE(\bbeta)
  =
  (\by-\bX\bbeta)^\top(\by-\bX\bbeta).
\]
对 $\bbeta$ 求梯度并令其为零：
\[
  -2\bX^\top(\by-\bX\hat{\bbeta})=0.
\]
得到正规方程
\[
  \boxed{\bX^\top\bX\hat{\bbeta}=\bX^\top\by.}
\]
若 $\bX$ 列满秩，使 $\bX^\top\bX$ 可逆，则
\[
  \boxed{
  \hat{\bbeta}
  =
  (\bX^\top\bX)^{-1}\bX^\top\by.
  }
\]

\begin{methodbox}[多元回归理论证明题]
教师录音确认：multiple regression 不考给出大数据集后的繁重 numerical calculation。应重点准备以下证明路线：
\begin{enumerate}
  \item 从最小化 $SSE(\bbeta)$ 推出正规方程 $\bX^\top\bX\hat{\bbeta}=\bX^\top\by$；
  \item 在列满秩条件下写出 $\hat{\bbeta}=(\bX^\top\bX)^{-1}\bX^\top\by$；
  \item 由 $\by=\bX\bbeta+\bvarepsilon$ 推出 $\E(\hat{\bbeta})=\bbeta$；
  \item 用协方差矩阵运算推出 $\Var(\hat{\bbeta})=\sigma^2(\bX^\top\bX)^{-1}$；
  \item 用 $\be=(\bI-\bH)\by$ 与 $\bX^\top\be=\mathbf0$ 证明残差性质。
\end{enumerate}
\end{methodbox}

\subsection{Hat matrix 与残差矩阵}

定义帽子矩阵
\[
  \boxed{
  \bH=\bX(\bX^\top\bX)^{-1}\bX^\top.
  }
\]
则
\[
  \hat{\by}=\bH\by,
  \qquad
  \be=\by-\hat{\by}=(\bI-\bH)\by.
\]
$\bH$ 对称且幂等：
\[
  \bH^\top=\bH,
  \qquad
  \bH^2=\bH.
\]
残差与设计矩阵的每一列正交：
\[
  \boxed{\bX^\top\be=\mathbf 0.}
\]
若模型含截距列 $\1$，则 $\1^\top\be=0$，即残差和为零。

\subsection{估计量与残差的分布性质}

若
\[
  \E(\bvarepsilon)=\mathbf0,
  \qquad
  \Var(\bvarepsilon)=\sigma^2\bI,
\]
则
\[
  \boxed{\E(\hat{\bbeta})=\bbeta,}
\]
\[
  \boxed{
  \Var(\hat{\bbeta})
  =
  \sigma^2(\bX^\top\bX)^{-1}.
  }
\]
证明：
\[
  \hat{\bbeta}
  =
  (\bX^\top\bX)^{-1}\bX^\top(\bX\bbeta+\bvarepsilon)
  =
  \bbeta+(\bX^\top\bX)^{-1}\bX^\top\bvarepsilon.
\]
取期望即得无偏性；利用协方差矩阵运算即得方差。

残差协方差矩阵：
\[
  \boxed{
  \Var(\be)=\sigma^2(\bI-\bH).
  }
\]
因此
\[
  \Var(e_i)=\sigma^2(1-h_{ii}),
\qquad
  \Cov(e_i,e_j)=-\sigma^2h_{ij}\quad(i\ne j).
\]
残差一般并不相互独立。

误差方差的无偏估计：
\[
  \boxed{
  s^2=\frac{SSE}{n-p},
  }
\]
其中 $p$ 是估计参数个数，\textbf{包括截距}。

% ============================================================
\section{录音确认的回归范围调整}
% ============================================================

\begin{scopebox}[Polynomial regression 不纳入期末]
教师录音中明确说明：polynomial regression 属于 computer-based / R-based / graph-based work，手算负担过重，因此不进入期末考试。此前根据黑板照片临时标注为重点的判断，以本次录音确认结果为准。

回归复习时间应转移到：
\begin{itemize}
  \item simple linear regression 的手算、最小二乘与基础推导；
  \item multiple linear regression 的矩阵最小二乘、估计量期望/方差/协方差与残差性质证明；
  \item Week 9 中 regression inference 的证明路线。
\end{itemize}
\end{scopebox}

% ============================================================
\section{MLE：从最大似然到推断}
% ============================================================

\subsection{Likelihood、log-likelihood 与 score}

设 $Y_1,\ldots,Y_n$ 独立同分布，密度或概率质量函数为
$f(y;\theta)$。观测值为 $y_1,\ldots,y_n$。
\[
  \boxed{
  L(\theta)
  =
  \prod_{i=1}^nf(y_i;\theta).
  }
\]
\[
  \boxed{
  \ell(\theta)
  =
  \log L(\theta)
  =
  \sum_{i=1}^n\log f(y_i;\theta).
  }
\]
score 函数为
\[
  \boxed{
  U(\theta)
  =
  \frac{\partial\ell(\theta)}{\partial\theta}.
  }
\]
MLE 定义为
\[
  \boxed{
  \hat\theta_{\mathrm{MLE}}
  =
  \arg\max_{\theta\in\Theta}L(\theta)
  =
  \arg\max_{\theta\in\Theta}\ell(\theta).
  }
\]

\begin{methodbox}[MLE 推导标准流程]
\begin{enumerate}
  \item 写出单个观测的 pmf/pdf，并标明参数空间；
  \item 利用独立性写乘积似然；
  \item 取对数，将乘积转为求和；
  \item 求 score，令 $U(\theta)=0$；
  \item 解出候选 MLE；
  \item 用二阶导数、凹性、端点或参数空间检查最大值；
  \item 计算 Fisher 信息，写出大样本标准误与近似推断。
\end{enumerate}
\end{methodbox}

\subsection{Fisher 信息与大样本推断}

对标量参数，
\[
  \boxed{
  I_n(\theta)
  =
  -\E_\theta\left[
  \frac{\partial^2\ell(\theta)}{\partial\theta^2}
  \right].
  }
\]
在常见正则条件下，也有
\[
  I_n(\theta)=\Var_\theta[U(\theta)].
\]
独立观测的信息相加：
\[
  I_n(\theta)=nI_1(\theta).
\]

\begin{memorybox}[MLE 的渐近正态性]
在常见正则条件下，当 $n$ 大时，
\[
  \boxed{
  \hat\theta_{\mathrm{MLE}}
  \approx
  N\left(\theta,\ I_n(\theta)^{-1}\right).
  }
\]
用 MLE 代入未知参数后，
\[
  \boxed{
  \widehat{\SE}(\hat\theta_{\mathrm{MLE}})
  =
  \frac{1}{\sqrt{I_n(\hat\theta_{\mathrm{MLE}})}}.
  }
\]
近似 $(1-\alpha)100\%$ Wald CI：
\[
  \boxed{
  \hat\theta_{\mathrm{MLE}}
  \pm
  z_{\alpha/2}
  \widehat{\SE}(\hat\theta_{\mathrm{MLE}}).
  }
\]
\end{memorybox}

\begin{warnbox}[Wald 近似的适用边界]
Wald CI 与 Wald 检验依赖大样本渐近正态近似。样本很小、似然明显偏斜，或参数接近边界时，近似可能较差；区间甚至可能越出合法参数空间，例如比例区间超出 $[0,1]$、正参数区间出现负端点。遇到此类结果，应明确说明这是简单 Wald 近似的局限。
\end{warnbox}

检验 $H_0:\theta=\theta_0$ 可使用近似 Wald 统计量
\[
  \boxed{
  Z_{\mathrm{obs}}
  =
  \frac{\hat\theta_{\mathrm{MLE}}-\theta_0}
  {1/\sqrt{I_n(\hat\theta_{\mathrm{MLE}})}}.
  }
\]
再按单尾或双尾临界值法作决策。

\begin{conceptbox}[Fisher 信息的直觉]
信息越大，对数似然在真参数附近越尖，参数越容易被精确定位；因此估计方差约为信息的倒数。样本量增加时信息通常线性增加，标准误通常按 $1/\sqrt n$ 缩小。
\end{conceptbox}

\subsection{MLE 的不变性}

若 $\hat\theta$ 是 $\theta$ 的 MLE，且 $\phi=g(\theta)$，则在适当条件下
\[
  \boxed{
  \hat\phi_{\mathrm{MLE}}=g(\hat\theta_{\mathrm{MLE}}).
  }
\]
例如指数分布 rate 参数为 $\lambda$，mean 参数为
$\beta=1/\lambda$；若 $\hat\lambda=1/\bar Y$，则
\[
  \hat\beta=\frac1{\hat\lambda}=\bar Y.
\]

% ============================================================
\section{四类分布的 MLE 与推断}
% ============================================================

\subsection{正态分布}

\subsubsection{方差已知，均值未知}

若
\[
  Y_1,\ldots,Y_n\iid N(\mu,\sigma^2),
\]
且 $\sigma^2$ 已知，则
\[
  \ell(\mu)
  =
  -\frac n2\log(2\pi\sigma^2)
  -
  \frac1{2\sigma^2}\sum_{i=1}^n(y_i-\mu)^2.
\]
\[
  U(\mu)
  =
  \frac1{\sigma^2}\sum_{i=1}^n(y_i-\mu).
\]
令 $U(\mu)=0$：
\[
  \boxed{\hat\mu_{\mathrm{MLE}}=\bar y.}
\]
二阶导数为 $-n/\sigma^2<0$，所以确为最大值。
\[
  \boxed{
  I_n(\mu)=\frac n{\sigma^2},
  \qquad
  \SE(\hat\mu)=\frac{\sigma}{\sqrt n}.
  }
\]
此处正态均值 MLE 的正态分布是精确结果，而不仅是渐近结果。

\subsubsection{均值已知，方差未知}

令 $v=\sigma^2$，且 $\mu$ 已知：
\[
  \ell(v)
  =
  -\frac n2\log(2\pi)
  -\frac n2\log v
  -\frac1{2v}\sum_{i=1}^n(y_i-\mu)^2.
\]
求导并令其为零：
\[
  \boxed{
  \hat v_{\mathrm{MLE}}
  =
  \frac1n\sum_{i=1}^n(y_i-\mu)^2.
  }
\]
Fisher 信息：
\[
  \boxed{
  I_n(v)=\frac{n}{2v^2},
  \qquad
  \widehat{\SE}(\hat v)
  =
  \sqrt{\frac{2\hat v^2}{n}}
  =
  \hat v\sqrt{\frac2n}.
  }
\]

\subsubsection{均值与方差都未知}

\[
  \boxed{
  \hat\mu_{\mathrm{MLE}}=\bar y,
  \qquad
  \hat\sigma_{\mathrm{MLE}}^2
  =
  \frac1n\sum_{i=1}^n(y_i-\bar y)^2.
  }
\]
注意
\[
  \E(\hat\sigma_{\mathrm{MLE}}^2)
  =
  \frac{n-1}{n}\sigma^2,
\]
所以它有偏；无偏样本方差使用分母 $n-1$。

令 $v=\sigma^2$，Fisher 信息矩阵为
\[
  \boxed{
  \mathbf I_n(\mu,v)
  =
  \begin{pmatrix}
    n/v & 0\\
    0 & n/(2v^2)
  \end{pmatrix}.
  }
\]
交叉信息为零，因为相关交叉二阶导数的期望包含
\[
  \E\left[\sum(Y_i-\mu)\right]=0.
\]
因此 $\hat\mu$ 与 $\hat v$ 在大样本下渐近不相关。

\subsection{二项分布}

先考虑一个二项计数
\[
  Y\sim\operatorname{Binomial}(m,p),
\]
观测到 $y$ 次成功。似然与对数似然为
\[
  L(p)
  =
  \binom myp^y(1-p)^{m-y},
\]
\[
  \ell(p)
  =
  \log\binom my+y\log p+(m-y)\log(1-p).
\]
score：
\[
  U(p)=\frac yp-\frac{m-y}{1-p}.
\]
当 $0<y<m$ 时，令其为零：
\[
  \boxed{\hat p_{\mathrm{MLE}}=\frac ym.}
\]
若 $y=0$ 或 $y=m$，score 方程在参数空间内部没有根，但检查边界仍得到同一公式：
\[
  y=0\Rightarrow\hat p=0,
  \qquad
  y=m\Rightarrow\hat p=1.
\]
Fisher 信息：
\[
  \boxed{
  I(p)=\frac{m}{p(1-p)}.
  }
\]
所以
\[
  \boxed{
  \widehat{\SE}(\hat p)
  =
  \sqrt{\frac{\hat p(1-\hat p)}{m}}.
  }
\]

若有 $n$ 个独立的 $\operatorname{Binomial}(m,p)$ 观测，则
\[
  \hat p=\frac{\sum_{i=1}^ny_i}{nm},
  \qquad
  I_n(p)=\frac{nm}{p(1-p)}.
\]
Bernoulli 是 $m=1$ 的特殊情况。

\begin{examplebox}[二项 MLE 与 Fisher 信息]
若 $m=50$ 次试验观察到 $y=18$ 次成功，则
\[
  \hat p=\frac{18}{50}=0.36,
\]
\[
  \widehat{\SE}(\hat p)
  =
  \sqrt{\frac{0.36(0.64)}{50}}
  \approx0.0679.
\]
\end{examplebox}

\subsection{Poisson 分布}

若
\[
  Y_1,\ldots,Y_n\iid\operatorname{Poisson}(\lambda),
\]
则
\[
  f(y_i;\lambda)
  =
  \frac{\lambda^{y_i}e^{-\lambda}}{y_i!}.
\]
对数似然：
\[
  \ell(\lambda)
  =
  \sum_{i=1}^ny_i\log\lambda
  -
  n\lambda
  -
  \sum_{i=1}^n\log(y_i!).
\]
score：
\[
  U(\lambda)
  =
  \frac{\sum y_i}{\lambda}-n.
\]
因此
\[
  \boxed{\hat\lambda_{\mathrm{MLE}}=\bar y.}
\]
二阶导数：
\[
  \ell''(\lambda)
  =
  -\frac{\sum y_i}{\lambda^2}.
\]
利用 $\E(\sum Y_i)=n\lambda$：
\[
  \boxed{
  I_n(\lambda)=\frac n\lambda.
  }
\]
因此
\[
  \boxed{
  \widehat{\SE}(\hat\lambda)
  =
  \sqrt{\frac{\hat\lambda}{n}}.
  }
\]
近似 CI：
\[
  \boxed{
  \hat\lambda
  \pm
  z_{\alpha/2}\sqrt{\frac{\hat\lambda}{n}}.
  }
\]

\subsection{指数分布}

课程采用 rate 参数化：
\[
  Y_1,\ldots,Y_n\iid\operatorname{Exponential}(\theta),
\qquad
  f(y;\theta)=\theta e^{-\theta y},
\quad y>0,\theta>0.
\]
此时
\[
  \E(Y)=\frac1\theta,
  \qquad
  \Var(Y)=\frac1{\theta^2}.
\]
对数似然：
\[
  \ell(\theta)
  =
  n\log\theta-\theta\sum_{i=1}^ny_i.
\]
score：
\[
  U(\theta)
  =
  \frac n\theta-\sum_{i=1}^ny_i.
\]
因此
\[
  \boxed{
  \hat\theta_{\mathrm{MLE}}
  =
  \frac{n}{\sum y_i}
  =
  \frac1{\bar y}.
  }
\]
\[
  \ell''(\theta)=-\frac n{\theta^2},
\qquad
  \boxed{I_n(\theta)=\frac n{\theta^2}.}
\]
所以
\[
  \boxed{
  \widehat{\SE}(\hat\theta)
  =
  \frac{\hat\theta}{\sqrt n}.
  }
\]
近似 CI：
\[
  \boxed{
  \hat\theta
  \pm
  z_{\alpha/2}\frac{\hat\theta}{\sqrt n}.
  }
\]

\begin{warnbox}[指数分布参数化]
必须先确认参数是 rate 还是 mean/scale。
\[
  f(y;\theta)=\theta e^{-\theta y}
  \quad\Rightarrow\quad
  \theta\text{ 是 rate，MLE}=1/\bar y.
\]
\[
  f(y;\beta)=\frac1\beta e^{-y/\beta}
  \quad\Rightarrow\quad
  \beta\text{ 是 mean/scale，MLE}=\bar y.
\]
\end{warnbox}

\subsection{四类分布速查表}

\begin{center}
\renewcommand{\arraystretch}{1.35}
\small
\begin{tabularx}{\textwidth}{p{2.5cm}p{3.1cm}p{3.1cm}X}
\toprule
模型 & MLE & Fisher 信息 & 估计标准误\\
\midrule
$N(\mu,\sigma^2)$，$\sigma^2$ 已知
& $\hat\mu=\bar Y$
& $I_n(\mu)=n/\sigma^2$
& $\sigma/\sqrt n$\\
$N(\mu,v)$，$\mu$ 已知
& $\hat v=n^{-1}\sum(Y_i-\mu)^2$
& $I_n(v)=n/(2v^2)$
& $\hat v\sqrt{2/n}$\\
$\operatorname{Binomial}(m,p)$
& $\hat p=Y/m$
& $I(p)=m/[p(1-p)]$
& $\sqrt{\hat p(1-\hat p)/m}$\\
$\operatorname{Poisson}(\lambda)$
& $\hat\lambda=\bar Y$
& $I_n(\lambda)=n/\lambda$
& $\sqrt{\hat\lambda/n}$\\
$\operatorname{Exponential}(\theta)$，rate
& $\hat\theta=1/\bar Y$
& $I_n(\theta)=n/\theta^2$
& $\hat\theta/\sqrt n$\\
\bottomrule
\end{tabularx}
\end{center}

% ============================================================
\section{高频证明与答题模板}
% ============================================================

\subsection{证明题的通用路线}

\begin{methodbox}[证明估计量无偏]
\begin{enumerate}
  \item 把估计量写成随机变量的线性组合；
  \item 利用期望线性性；
  \item 代入每个随机变量的期望；
  \item 化简为目标参数。
\end{enumerate}
\end{methodbox}

\begin{methodbox}[计算估计量方差]
\begin{enumerate}
  \item 提取常数的平方；
  \item 展开线性组合方差；
  \item 使用独立性删除协方差项，或保留给定协方差；
  \item 化简并检查单位。
\end{enumerate}
\end{methodbox}

\begin{methodbox}[证明 OLS 残差性质]
直接从正规方程出发：
\[
  \bX^\top(\by-\bX\hat{\bbeta})=\mathbf0
  \quad\Longrightarrow\quad
  \bX^\top\be=\mathbf0.
\]
若 $\bX$ 第一列是全 1，则第一条方程给出 $\sum e_i=0$；其他列分别给出残差与各预测变量正交。
\end{methodbox}

\subsection{完整书面结论的语言}

\begin{itemize}
  \item \textbf{拒绝 $H_0$：}“在 $\alpha=0.05$ 显著性水平下，数据提供充分证据支持……”
  \item \textbf{不能拒绝 $H_0$：}“在 $\alpha=0.05$ 显著性水平下，数据不足以支持……” 
  \item 不写“证明了 $H_0$ 为真”，也不写“接受 $H_0$”。
  \item CI 解释必须包含参数、置信水平、方向与单位。
  \item 回归斜率解释必须包含“保持其他变量不变”（多元回归）与响应单位。
\end{itemize}

% ============================================================
\section{考前公式总表}
% ============================================================

\subsection{无偏估计与 MSE}

\[
  \Bias(\hat\theta)=\E(\hat\theta)-\theta,
  \qquad
  \MSE(\hat\theta)=\Var(\hat\theta)+\Bias(\hat\theta)^2.
\]
\[
  a^*_{\text{combine}}
  =
  \frac{v_2-c}{v_1+v_2-2c}.
\]

\subsection{均值推断}

\[
  \bar x\pm z_{\alpha/2}\frac{\sigma}{\sqrt n},
  \qquad
  Z=\frac{\bar x-\mu_0}{\sigma/\sqrt n}.
\]
\[
  \bar x\pm t_{\alpha/2,n-1}\frac{s}{\sqrt n},
  \qquad
  t=\frac{\bar x-\mu_0}{s/\sqrt n}.
\]
\[
  (\bar x_1-\bar x_2)
  \pm
  z_{\alpha/2}
  \sqrt{\frac{s_1^2}{n_1}+\frac{s_2^2}{n_2}}.
\]
\[
  s_p^2
  =
  \frac{(n_1-1)s_1^2+(n_2-1)s_2^2}{n_1+n_2-2},
\]
\[
  t
  =
  \frac{(\bar x_1-\bar x_2)-\Delta_0}
  {s_p\sqrt{1/n_1+1/n_2}}.
\]
\[
  t_{\mathrm{paired}}
  =
  \frac{\bar d-\Delta_0}{s_d/\sqrt n}.
\]

\subsection{比例推断}

\[
  \hat p=\frac xn,
  \qquad
  \hat p\pm z_{\alpha/2}
  \sqrt{\frac{\hat p(1-\hat p)}n},
\]
\[
  Z=\frac{\hat p-p_0}{\sqrt{p_0(1-p_0)/n}}.
\]
\[
  (\hat p_1-\hat p_2)
  \pm
  z_{\alpha/2}
  \sqrt{
  \frac{\hat p_1(1-\hat p_1)}{n_1}
  +
  \frac{\hat p_2(1-\hat p_2)}{n_2}
  }.
\]
\[
  \hat p_{\mathrm{pool}}=\frac{x_1+x_2}{n_1+n_2},
\quad
  Z=
  \frac{\hat p_1-\hat p_2}
  {\sqrt{\hat p_{\mathrm{pool}}(1-\hat p_{\mathrm{pool}})
  (1/n_1+1/n_2)}}.
\]

\subsection{回归}

\[
  \hat\beta_1=\frac{S_{xy}}{S_{xx}},
  \qquad
  \hat\beta_0=\bar y-\hat\beta_1\bar x,
\]
\[
  SSE=S_{yy}-\frac{S_{xy}^2}{S_{xx}},
  \qquad
  s_e^2=\frac{SSE}{n-2}.
\]
\[
  \Var(\hat\beta_1)=\frac{\sigma^2}{S_{xx}},
  \qquad
  \Var(\hat\beta_0)
  =
  \sigma^2\left(\frac1n+\frac{\bar x^2}{S_{xx}}\right),
\]
\[
  \Cov(\hat\beta_0,\hat\beta_1)
  =
  -\frac{\sigma^2\bar x}{S_{xx}}.
\]
\[
  \hat{\bbeta}=(\bX^\top\bX)^{-1}\bX^\top\by,
  \qquad
  \bX^\top\be=\mathbf0,
\]
\[
  \Var(\hat{\bbeta})
  =
  \sigma^2(\bX^\top\bX)^{-1},
  \qquad
  s^2=\frac{SSE}{n-p}.
\]
\[
  \frac{d}{dx}
  \left(\beta_0+\beta_1x+\cdots+\beta_qx^q\right)
  =
  \beta_1+2\beta_2x+\cdots+q\beta_qx^{q-1}.
\]
\[
  \text{二次模型转折点：}\qquad
  x^\star=-\frac{\beta_1}{2\beta_2}.
\]
\[
  z=x-c:\qquad
  \gamma_0=m(c),\quad
  \gamma_1=m'(c),\quad
  \gamma_2=\beta_2.
\]

\subsection{MLE}

\[
  L(\theta)=\prod_{i=1}^nf(y_i;\theta),
  \quad
  \ell(\theta)=\log L(\theta),
  \quad
  U(\theta)=\ell'(\theta),
\]
\[
  I_n(\theta)=-\E[\ell''(\theta)],
  \qquad
  \widehat{\SE}(\hat\theta)
  =
  \frac1{\sqrt{I_n(\hat\theta)}}.
\]
\[
  \hat\theta
  \pm
  z_{\alpha/2}\widehat{\SE}(\hat\theta).
\]

% ============================================================
\section{最后检查清单}
% ============================================================

\begin{scopebox}[交卷前逐项检查]
\begin{enumerate}
  \item 我是否先定义了参数，而不只是写样本统计量？
  \item 单尾还是双尾是否由题意决定，临界值是否放在正确尾部？
  \item 未知方差时是否正确使用 $t$ 与自由度？
  \item 两样本是独立还是配对？等方差是否有依据？
  \item 比例 CI 与比例检验是否使用了正确的标准误？
  \item 回归设计矩阵是否含截距列？参数个数 $p$ 是否包括截距？
  \item 多元回归系数解释是否写了“保持其他变量不变”？
  \item 多元回归证明是否写清 $\bX$ 列满秩、含截距、参数个数 $p$ 与残差自由度 $n-p$？
  \item MLE 是否写了参数空间、score、边界或最大值检查与 Fisher 信息？
  \item 最终结论是否联系题目语境与单位，并使用“拒绝/不能拒绝 $H_0$”？
\end{enumerate}
\end{scopebox}

\begin{scopebox}[资料核对与暂留问题]
本讲义已使用以下材料交叉核对：课程中文讲义、英文讲义、Week 1--14 相关课件、两套 2026 样卷及答案、往年试题中与当前范围相符的部分。

第三张黑板图中额外补写统计量的具体考试含义，等待教师后续确认。确认后应优先核对：
\begin{itemize}
  \item 它是否只是单样本/双样本未知方差 $t$ 统计量的提示；
  \item 是否要求额外的总体方差本身推断；
  \item 是否给定等方差条件，或要求使用其他双样本标准误。
\end{itemize}
\end{scopebox}

\end{document}
